\documentclass[11pt]{article}
\usepackage[T1]{fontenc}
\usepackage[margin=1in]{geometry}
\usepackage{microtype}
\usepackage{xcolor}
\usepackage[colorlinks=true,linkcolor=blue,urlcolor=blue]{hyperref}
\usepackage{xurl}
\usepackage{enumitem}

\setlength{\parindent}{0pt}
\setlength{\parskip}{6pt}
\definecolor{responseblue}{RGB}{24,72,120}

\newcommand{\reviewcomment}[2]{%
  \section*{Reviewer Comment #1}
  \begin{quote}\itshape #2\end{quote}}
\newcommand{\response}{\subsection*{Response}}
\newcommand{\manuscriptchanges}{\subsection*{Changes in the manuscript}}
\newcommand{\computationalchanges}{\subsection*{Computational/reproducibility changes}}
\newcommand{\status}[1]{\subsection*{Status}\textbf{#1}}

\begin{document}

\begin{center}
{\Large\bfseries Response to Reviewer}\par
\vspace{4pt}
{\large Manuscript JBCB-1505}\par
\vspace{2pt}
\textit{A formal and reproducible framework for auditing observable behavior in
qualitative biological network models}
\end{center}

We thank the reviewer for a careful and constructive report. We particularly
appreciate the positive assessment of the formal definitions, fixed-point
algorithm, mCRL2 comparison, independent simulation game, and reproducibility
materials. We agree with the central criticism: the submitted version
established computational correctness more convincingly than biological
utility. The revision therefore does more than add exposition. It introduces a
three-interface GIM experiment on unchanged Petri-net structures, reorganizes
the applications around an explicit DNA-damage question, adds verified
biological support for the interface design, removes the unjustified ordinal
HPN class analysis, redesigns the central figures, and separates four evidence
layers throughout. Negative HPN results remain unchanged and visible.

\reviewcomment{R1}{The manuscript demonstrates computational correctness much
more convincingly than biological utility; what important biological conclusion
can be obtained that existing comparison approaches would not reveal?}

\response
We agree. The revised central result is no longer the six constructed controls.
It is the interface-sensitivity analysis of the GIM animal/human and
\textit{Arabidopsis} DNA-damage models. The unchanged Petri-net structures have
PN-GDDA similarity 0.9975846. Nevertheless, the behavioral result depends on
the biological readout: Interfaces A and B are weakly bisimilar, while Interface
C is non-comparable in both simulation directions and has $d_6=0.3529$. Thus a
structural comparison reports nearly identical selected wiring, whereas the
behavioral comparison locates the observational resolution at which the encoded
correspondence breaks. This is the substantive model-level insight. We explicitly
avoid extending it to organism-level equivalence.

\manuscriptchanges
The Abstract and Introduction now lead with the biological problem. Sections
4.1--4.2 formulate the question, present independent biological background, and
report the structural/behavioral discrepancy. Table 3 and Fig. 3 report the
executed A/B/C result. Section 5.1 interprets structural and behavioral
information as complementary.

\computationalchanges
Added \nolinkurl{src/gim_interface_analysis.py},
\nolinkurl{results/gim_interface_sensitivity.csv}, and regression tests in
\nolinkurl{tests/test_gim_interface_analysis.py}. The result is regenerated by
\texttt{make gim-interface} and \texttt{make figures}.

\status{Addressed by reframing and additional analysis}

\reviewcomment{R2}{The biological part should be organized around one or two
explicit biological questions, moving from structural comparison to behavioral
comparison, biological meaning, independent evidence, and a testable
hypothesis.}

\response
We agree and reorganized the applications accordingly. The primary question is:
at what observational resolution do the two encoded DNA-damage control
structures preserve common behavior, and when do mechanism-specific outcomes
make that correspondence fail? The sequence is now biological background,
structural result, coarse and intermediate behavioral results, mechanism-resolved
result, interpretation, and a limited model-derived hypothesis. RCD is retained
as a secondary contrasting question about directional containment, rather than
being given equal evidential weight.

\manuscriptchanges
Sections 4.1--4.2 implement the requested sequence. The concluding GIM paragraph
states that readouts aggregating terminal loss of proliferative capacity are
expected to be more compatible with A/B, whereas readouts distinguishing
apoptosis, SMR induction, and differentiation/endoreduplication should break the
apparent relation. This is labeled a testable model hypothesis, not an
experimental finding. Section 4.5 presents RCD as secondary.

\computationalchanges
The notebook now executes and displays the three GIM interfaces before the HPN
analysis. Interface rationale, labels, internal transitions, state counts,
relations, diagnostics, interpretations, and literature keys are exported to
CSV.

\status{Fully addressed}

\reviewcomment{R3}{The GIM example is promising but needs a deeper explanation
of why the comparison is biologically interesting, which functions are shared,
which mechanisms differ, and what follows from weak bisimilarity.}

\response
We agree. The revision distinguishes shared high-level functions from distinct
implementations. Both encodings contain damage sensing, checkpoint control,
repair, and recovery. The animal/human encoding includes CHK/p53-related control
and apoptosis, whereas the plant encoding includes the plant-specific SOG1,
WEE1 and SMR program and a combined differentiation/endoreduplication outcome.
The text does not describe p53 and SOG1 as orthologs and does not equate plant
programmed cell death with mammalian apoptosis. Weak bisimilarity is stated only
with respect to the declared A/B observations.

\manuscriptchanges
The motivating biological background appears in the Introduction and Section
4.1. Table 2 provides process-by-process interpretation and literature support.
Section 4.2 explains the matched paths and the limits of the result. The full
executable nets remain available at readable scale in Appendix A (Figs. 5--6).

\computationalchanges
The plant transition previously named \nolinkurl{CHK_signal} was renamed
\nolinkurl{SOG1_program} without changing topology or reachability. Interface C
exposes separate \nolinkurl{apoptosis}, \nolinkurl{smr_induction}, and
\nolinkurl{differentiation_or_endoreduplication} labels.

\status{Addressed by reframing and additional analysis}

\reviewcomment{R4}{The observational interface is a biological modeling
assumption. The manuscript should justify observable and silent events and show
how conclusions change under several biologically plausible interfaces.}

\response
We agree completely. Three interfaces are now defined and executed on unchanged
structures. A observes coarse functions. B distinguishes DSB and
replication-stress detection, ATM/ATR-associated signaling, and broad repair
families while collapsing terminal fate. C additionally distinguishes the
organism-specific terminal events. The outcome was not forced: A and B are
weakly bisimilar; C loses weak bisimilarity and both simulations. The structural
score is constant, which isolates interface resolution as the changed factor.

\manuscriptchanges
Section 4.1 defines A/B/C and Table 2 justifies each observable/internal choice.
Table 3 reports both simulation directions and Fig. 3 visualizes the interface
refinement. Section 5.2 states explicitly that $\Sigma$ is a biological
hypothesis.

\computationalchanges
Interfaces and literature keys are versioned in
\nolinkurl{src/concurrent_biomodels.py}. The files
\nolinkurl{gim_interface_sensitivity.csv},
\nolinkurl{gim_interface_justification.csv}, and a generated LaTeX table are
written deterministically. Tests assert unchanged transition sets, 13/14 states,
constant PN-GDDA, and the exact A/B/C classes.

\status{Fully addressed}

\reviewcomment{R5}{The six construction-level cases verify implementation
behavior but are not biological validation and should not be presented as such.}

\response
We agree. The six cases are now consistently called a controlled
construction-level software/formal benchmark. Their 6/6 result verifies that the
implementation recovers the declared formal distinctions; it is not evidence of
biological usefulness.

\manuscriptchanges
Section 3 and the caption of Fig. 2 state both what the benchmark establishes and
what it does not. The Abstract, Introduction, Discussion, Conclusion, README,
and reproducibility guide use the same terminology.

\computationalchanges
The benchmark code and expected outcomes are unchanged. Figure-generation text
and documentation were revised so the output is labeled ``software verification
only.''

\status{Fully addressed}

\reviewcomment{R6}{The comparison with PN-GDDA should demonstrate
complementarity rather than generic superiority, preferably on a real
biological example.}

\response
We agree. PN-GDDA does not attempt to represent labels, initial markings, or
reachable branching, so its lower agreement on a behavioral construction task
cannot establish generic inferiority. The revised GIM experiment supplies the
requested biological demonstration: PN-GDDA remains 0.9976 at all interfaces,
whereas the behavioral relation changes at Interface C. The methods therefore
answer complementary questions.

\manuscriptchanges
The Fig. 2 caption rejects a superiority interpretation. Sections 4.2, 5.1, and
the Conclusion explain the GIM structural/behavioral discrepancy and avoid
``outperforms'' or ``better'' claims.

\computationalchanges
Both 592-slot and 576-orbit PN-GDDA values remain in the machine-readable output.
No threshold or score was changed to improve the result.

\status{Fully addressed}

\reviewcomment{R7}{The HPN-DREAM/CASPOTS analysis is too limited for strong
biological conclusions and should be presented as exploratory or proof of
concept; its negative results should be retained.}

\response
We agree. All nine pair--condition observations and all negative results are
retained. The trace association remains $\rho=-0.596$ with exact $p=0.556$; the
LTS-GDA distance association remains $\rho=0.669$ with $p=0.097$; and the
CASPOTS native-context advantage remains unestablished at exact $p=0.25$. We
describe this analysis as an exploratory held-out compatibility check, not as
validation of predictive utility, prognosis, or cell-line specificity.

\manuscriptchanges
Section 4.4 and Fig. 4 are explicitly labeled exploratory. Section 5.3 gives the
sample-size and metric limitations, and the Abstract and Conclusion preserve the
negative interpretation.

\computationalchanges
No favorable alternative test was sought. The existing exact permutations for
the two predeclared continuous diagnostics and the exact CASPOTS sign test were
retained.

\status{Addressed within the limitations of the available data}

\reviewcomment{R8}{The formal relations do not form an obvious one-dimensional
ordinal scale, especially because one-way simulation is directional. The
Spearman encoding risks discarding the information that motivates simulation.}

\response
We agree and removed the ordinal analysis rather than defending it. Strong/weak
equivalence, directional simulation, and non-comparability do not have a
biologically justified total order. The revised output preserves three
left-in-right relations, three right-in-left relations, and three pairs with no
simulation. Because the experimental pair RMSE is symmetric and $n=9$, we make
no direction-aware inferential claim. Descriptive medians are shown without a
scalar class test.

\manuscriptchanges
The former $\rho=0.091$, $p=0.111$ result is removed from the manuscript and all
public documentation because the analysis itself was invalid, not because it
was negative. Section 4.4 explains the decision; Fig. 4b displays the nominal
directions; Section 5.3 discusses the limitation.

\computationalchanges
\nolinkurl{src/hpn_dream_validation.py} no longer creates ordinal class scores or
formal-class Spearman fields. It writes \nolinkurl{relation_direction} and
nominal summaries. Regression tests require 3/3/3 directions and explicitly
reject the obsolete fields. The result JSON records
\nolinkurl{formal_class_scalar_encoding_used=false} and
\nolinkurl{formal_class_inferential_test_performed=false}.

\status{Fully addressed}

\reviewcomment{R9}{Machine-checked transition provenance establishes
reproducibility of modeling choices, not that the resulting curated Petri nets
are biologically correct.}

\response
We agree. The revision distinguishes provenance, internal reproducibility,
executable-model validation, and external validation of a biological
interpretation. The curated models are described as literature-informed
encodings without kinetics, concentrations, tissue context, stochasticity, or
independent experimental calibration.

\manuscriptchanges
Section 4.5 states what transition-level provenance does and does not establish.
Section 5.3 and Table 7 separate the four evidence layers. GIM conclusions are
labeled as model-derived, literature-consistent, or testable rather than
experimentally validated.

\computationalchanges
Provenance coverage tests are retained, but documentation no longer promotes
coverage to biological validation. The GIM justification CSV separately records
biological rationale and literature keys.

\status{Fully addressed}

\reviewcomment{R10}{The benchmark and central biological figures are difficult
to read at manuscript scale and require conceptual redesign rather than a higher
resolution export.}

\response
We agree. Both figures were regenerated as vector artwork with larger type and
less visual density. The benchmark is split into a decision matrix and a simple
agreement panel, explicitly labeled as software verification. The central GIM
figure no longer asks readers to decode two compressed Petri nets: it shows
shared functions, organism-specific implementations, matched coarse control
flow, and the A/B/C classification. Full nets are retained in the appendix at
readable width.

\manuscriptchanges
Revised Figs. 2--3 and their captions replace the submitted figures. Fig. 4 was
also redesigned to preserve HPN simulation direction instead of displaying an
ordinal class plot. Appendix Figs. 5--6 retain the detailed Petri nets.

\computationalchanges
All artwork is generated by \nolinkurl{make_figures.py}; the main scientific
figures are PDF vectors with PNG previews for inspection.

\status{Fully addressed}

\reviewcomment{R11}{The manuscript should separate formal correctness, software
implementation validation, validation of executable biological models, and
biological validation of conclusions.}

\response
We agree and adopted this four-layer taxonomy throughout. Formal proofs belong
to Layer I; construction controls and independent software oracles to Layer II;
hashes, ingestion controls, provenance, and semantics to Layer III; and
independent biological literature or held-out observations to Layer IV. Each
layer now has an explicit boundary.

\manuscriptchanges
The taxonomy is introduced in the Introduction and formalized in Discussion
Table 7. The Abstract, Section 3, Section 4, Discussion, and Conclusion use the
same distinctions.

\computationalchanges
README, \nolinkurl{REPRODUCIBILITY.md}, and the executed notebook now reproduce the
same evidence terminology and identify which artifact supports each layer.

\status{Fully addressed}

\reviewcomment{R12}{The manuscript currently reads as a formal-methods framework
illustrated with biological models; the revision should make one or two
carefully supported biological case studies stand out without weakening the
formal contribution.}

\response
We agree with this overall assessment. The formal framework, proofs, simulation
orientation, and external oracles are preserved. The narrative now begins with
the biological model-comparison problem, uses GIM as the central case, treats
RCD as a limited directional contrast, demotes transformed public models to
ingestion controls, and presents HPN as exploratory evidence. The conclusion is
now a biologically meaningful but restrained statement about locating the
observational boundary between high-level encoded function and
mechanism-specific behavior.

\manuscriptchanges
The Abstract, Introduction, all of Section 4, Discussion, and Conclusion were
rewritten. Supporting examples are compact, GIM is developed in depth, and all
organism-level or predictive overinterpretations are expressly excluded.

\computationalchanges
The revised manuscript, notebook, figures, deterministic result tables, tests,
and submission package are generated from the same repository state. mCRL2
202607.0 was rerun on the public, synthetic, and HPN LTSs.

\status{Addressed by reframing and additional analysis}

\section*{Summary of verification performed}

The complete notebook executed successfully. All 46 regression tests passed
with mCRL2 available. The revised JBCB source compiles without undefined
citations or references. The final source package is also compiled independently
from a flat directory, matching the journal upload constraint.

\end{document}
